% ============================================================
%  BAB 4 – PERANCANGAN DAN PENGEMBANGAN ARTEFAK
%  Compile standalone:  pdflatex bab-4.tex
%  Compile as part of thesis: pdflatex main.tex
% ============================================================
\documentclass[../../thesis]{subfiles}
\usepackage{hyphenat}
\usepackage{iftex}

% ── Encoding, Language & Fonts ───────────────────────────────
\ifXeTeX
    \usepackage{polyglossia}
    \setmainlanguage{bahasai}
    \usepackage{fontspec}
    \setmainfont{TeX Gyre Termes}
    \setsansfont{TeX Gyre Heros}
    \setmonofont[Scale=0.88]{TeX Gyre Cursor}
\else
    \usepackage[utf8]{inputenc}
    \usepackage[T1]{fontenc}
    \usepackage[indonesian]{babel}
    \usepackage{mathptmx}
    \usepackage{helvet}
    \usepackage{courier}
\fi

% ── Page geometry (A4, UI thesis margins) ────────────────────
\usepackage[
  a4paper,
  top    = 2.54cm,
  bottom = 2.54cm,
  left   = 3.17cm,
  right  = 2.54cm
]{geometry}

% ── Line spacing & paragraph ─────────────────────────────────
\usepackage{setspace}
\onehalfspacing
\setlength{\parindent}{1.27cm}
\setlength{\parskip}{0pt}

% ── Microtype & justification ────────────────────────────────
\usepackage{microtype}
\sloppy

% ── Headings ─────────────────────────────────────────────────
\usepackage{titlesec}

\titleformat{\chapter}[block]
  {\centering\rmfamily\bfseries\large}{}
  {0pt}{\MakeUppercase}
\titlespacing*{\chapter}{0pt}{0pt}{1.5em}

\titleformat{\section}
  {\rmfamily\bfseries\normalsize}{\thesection}{1em}{}
\titlespacing*{\section}{0pt}{1.2em}{0.6em}

\titleformat{\subsection}
  {\rmfamily\bfseries\normalsize}{\thesubsection}{1em}{}
\titlespacing*{\subsection}{0pt}{1em}{0.5em}

\titleformat{\subsubsection}
  {\rmfamily\bfseries\normalsize}{\thesubsubsection}{1em}{}
\titlespacing*{\subsubsection}{0pt}{0.8em}{0.4em}

% ── Tables ───────────────────────────────────────────────────
\usepackage{longtable, booktabs, array, multirow}
\usepackage{calc}
\usepackage{tabularx}

% ── Math ─────────────────────────────────────────────────────
\usepackage{amsmath, amssymb}

% ── Graphics ─────────────────────────────────────────────────
\usepackage{graphicx}
\makeatletter
\def\maxwidth{\ifdim\Gin@nat@width>\linewidth\linewidth\else\Gin@nat@width\fi}
\makeatother
\setkeys{Gin}{width=\maxwidth,keepaspectratio}

% ── Lists ────────────────────────────────────────────────────
\usepackage{enumitem}
\setlist[enumerate,1]{
  label      = \arabic*.,
  leftmargin = 1.27cm,
  itemsep    = 4pt,
  parsep     = 0pt,
  topsep     = 4pt
}
\setlist[enumerate,2]{
  label      = \alph*.,
  leftmargin = 1.27cm,
  itemsep    = 3pt,
  parsep     = 0pt,
  topsep     = 2pt
}

% ── Hyperlinks ───────────────────────────────────────────────
\usepackage[hidelinks]{hyperref}

% Bibliography setup is inherited from main.tex (biblatex + references.bib)
% when this file is compiled standalone via the subfiles class.

% ============================================================
\begin{document}

% When standalone: set counter so \chapter{} renders as "BAB 4".
% When in main.tex: main.tex sets the counter before \subfile{bab-4}.
\ifSubfilesClassLoaded{}{\setcounter{chapter}{3}}

\chapter{Perancangan dan Pengembangan Artefak}

Rancangan yang ditetapkan pada Bab~3 masih berupa \textit{blueprint} yang
nilainya baru teruji ketika diwujudkan menjadi sistem yang bekerja pada data
nyata. Bab ini
memaparkan perwujudan tersebut, yakni pembangunan artefak \textit{Knowledge
Graph Completion} (KGC) berbantuan \textit{Large Language Model} (LLM) tahap
demi tahap: dari mengubah buku teks PDF menjadi pengetahuan terstruktur,
menginstansiasikannya sebagai graf, hingga melengkapinya dengan relasi
lintas-disiplin yang tidak tersurat dalam satu sumber tunggal. Dipaparkan pula
arsitektur instrumen yang dikembangkan untuk mengukur artefak tersebut
(Subbab~\ref{sec:ingestion}--\ref{sec:instrumen-evaluasi}), sementara hasil
evaluasinya beserta pembahasan disajikan pada Bab~\ref{sec:hasil-evaluasi}.

% ============================================================
\section{Tahap Ingestion (PDF Parsing)}\label{sec:ingestion}
% ============================================================

Tahap ini mengubah \textit{e-book} PDF menjadi struktur data yang siap
diproses oleh LLM. Keluaran utama berupa teks per bab, metadata halaman,
daftar subbab, glosarium, dan materi pokok.

\begin{enumerate}
  \item \textbf{Ekstraksi teks PDF.} PDF dimuat menggunakan
        \texttt{SimpleDirectoryReader} (LlamaIndex) dengan \textit{backend}
        PyMuPDF, sehingga setiap halaman menjadi satu objek \texttt{Document}
        yang menyimpan \texttt{page\_label}.

  \item \textbf{Deteksi halaman TOC (daftar isi).} Halaman diklasifikasikan
        sebagai TOC bila memenuhi dua sinyal sekaligus: (a) memuat kata penanda
        seperti ``daftar isi'' atau ``bab'' (pencocokan tidak peka huruf
        besar--kecil), dan (b) memuat pola \textit{leader dots} yang diakhiri
        nomor halaman (contoh: ``..... 191''). Gabungan kedua sinyal ini
        mencegah \textit{false positive} pada halaman konten biasa yang kebetulan
        memuat hanya salah satu pola.

  \item \textbf{Parsing struktur Bab dan Subbab.} Pola Bab diekstrak dengan
        \textit{regex} yang menangkap format judul--\textit{leader dots}--nomor
        halaman (Contoh: ``Bab 4 Inovasi Bioteknologi..... 191''), sedangkan
        subbab diidentifikasi sebagai entri bertanda huruf besar (A., B., dst.)
        pada rentang TOC antara dua Bab berurutan. Hasil \textit{parsing}
        disusun sebagai daftar bab yang masing-masing memuat nama, halaman
        awal, dan daftar subbab.

  \item \textbf{Klasifikasi halaman TOC vs.\ konten.} Halaman yang terindikasi
        TOC pada deteksi awal digunakan untuk menyusun teks TOC gabungan,
        sementara halaman konten tidak disertakan agar \textit{parsing} bab dan
        subbab tetap stabil.

  \item \textbf{Ekstraksi glosarium.} Karena glosarium pada buku teks secara
        konvensional terletak di bagian akhir buku (\textit{back matter}),
        setelah seluruh materi pokok, pencarian penanda ``glosarium'' dibatasi
        hanya pada bagian setelah 30\% halaman pertama buku. Pembatasan ini
        menghindari \textit{false positive} ketika kata ``glosarium'' kebetulan
        muncul di bagian awal seperti daftar isi atau kata pengantar. Dari
        halaman penanda pertama yang memenuhi syarat, ekstraksi berlanjut halaman
        demi halaman hingga mencapai penanda awal bagian \textit{back matter}
        berikutnya, seperti ``Indeks'' atau ``Daftar Pustaka''. Pasangan
        istilah--definisi yang ditemukan diuraikan menjadi \textit{lookup table},
        yang kemudian dipakai sebagai validasi terminologi pada Tahap~2.

  \item \textbf{Penugasan bab ke \textit{chunk}.} Setiap \textit{chunk}
        dihubungkan ke bab berdasarkan interval halaman
        (\texttt{page\_start <= page < next\_page}), dengan \texttt{page\_label}
        dikonversi ke integer termasuk angka Romawi. Penugasan ini dilakukan satu
        kali untuk seluruh 788 \textit{chunk} dengan pencocokan rentang halaman
        tersebut, sehingga waktu komputasinya sangat kecil dibandingkan tahap
        ekstraksi berbasis LLM.

  \item \textbf{Ekstraksi materi pokok (Buku Guru).} Untuk tiap bab, sistem
        mencari \textit{region} teks yang sesuai judul bab di Buku Guru, lalu
        mengekstrak item di sekitar penanda ``Materi Pokok''. Hasilnya
        digunakan sebagai \textit{constraint} kurikulum pada Tahap~2 agar konsep
        yang diekstrak tetap sesuai kompetensi.
\end{enumerate}

% ============================================================
\section{Tahap Ekstraksi (LLM Prompting \& Schema)}\label{sec:ekstraksi}
% ============================================================

Tahap ini mengubah teks yang dihasilkan Tahap 1 menjadi \textit{knowledge
graph} terstruktur melalui \textit{schema-constrained prompting}.
Fokusnya adalah konsistensi struktur, validasi terminologi, dan
pembatasan konsep pada materi pokok.

\begin{enumerate}
  \item \textbf{Pemotongan teks (\textit{chunking}).}

  Teks dipecah menggunakan SentenceSplitter (LlamaIndex) dengan parameter
  sesuai Subbab~\ref{subsec:parameter-sistem}.

  \item \textbf{Agregasi per-bab.}

  \textit{Chunk} yang dipetakan ke bab yang sama digabungkan menjadi satu teks
  bab utuh, sehingga LLM memperoleh konteks bab secara penuh alih-alih potongan
  parsial per \textit{chunk}.

  \item \textbf{Ekstraksi triple.}

  Model: Gemini 2.5 Flash Lite via Google GenAI SDK. \textit{Prompt} bersifat
  \textit{few-shot} dan \textit{schema-constrained}. \textit{Prompt} menyertakan nama bab 
  beserta bab sebelum dan sesudahnya, daftar subbab dari daftar isi, glosarium sebagai 
  \textit{validation set}, materi pokok sebagai \textit{curriculum constraint}, konsep dari 
  buku lain sebagai \textit{cross-book context}, serta beberapa contoh \textit{triple} sebagai
  demonstrasi (\textit{few-shot}).

  \item \textbf{Struktur \textit{prompt}.}

  \textit{Prompt} didefinisikan \textit{inline} di dalam \textit{notebook} dan dibagi menjadi delapan blok:
  instruksi bahasa (keluaran wajib berbahasa Indonesia, kecuali istilah ilmiah
  tanpa padanan resmi); deklarasi hierarki graf (Kelas $\to$ Bab $\to$ Subbab
  $\to$ Konsep); filter konten (abaikan daftar isi, soal latihan, glosarium,
  rangkuman, dan elemen administratif); serta aturan ekstraksi konsep (delapan
  aturan), termasuk batas 3--5 konsep per subbab, deskripsi 1--2 kalimat ilmiah,
  dan validasi materi pokok Buku Guru. Templat \textit{prompt} ekstraksi
  selengkapnya (delapan blok dan empat blok dinamisnya) disajikan pada
  \ref{lamp:prompt-ekstraksi}.

  \item \textbf{Vokabulari relasi.}

  \textit{Prompt} membatasi ekstraksi relasi antar-konsep pada 15 tipe relasi
  intra-buku sesuai \textit{controlled vocabulary} ontologi
  (Subbab~\ref{subsec:skema-ontologi}); instruksi \textit{cross-chapter
  reasoning} secara terpisah menghasilkan relasi pedagogis antar-bab
  (\texttt{PRASYARAT} atau \texttt{MEMPERSIAPKAN}).

  \item \textbf{Penguraian keluaran terstruktur.}

  Respons LLM dibersihkan dari penanda \textit{markdown} lalu diuraikan dari
  format JSON menjadi objek terstruktur yang memuat ringkasan bab serta daftar
  konsep per subbab---lengkap dengan deskripsi, status validasi glosarium,
  rujukan materi pokok, tautan lintas-buku, dan relasi antar-konsep.

  \item \textbf{Validasi pascaproses.} Sebelum masuk ke graf, keluaran LLM
  dirapikan melalui tiga langkah. Dua langkah bertugas membuang keluaran yang
  tidak valid, sementara satu langkah justru mengisi keluaran apabila LLM gagal
  menghasilkannya:

  \begin{enumerate}
    \item \textit{Menyaring tautan lintas-buku.} Sebuah tautan dibuang jika buku
    tujuannya sama dengan buku asalnya, atau jika buku tujuannya tidak terdapat
    dalam korpus. Langkah ini mencegah munculnya tautan yang menunjuk ke dirinya
    sendiri (\textit{self-loop}) maupun tautan yang mengarah ke buku yang tidak
    tersedia (\textit{dangling}).

    \item \textit{Cadangan saat keluaran kosong.} Apabila LLM tidak menghasilkan
    satu pun tautan lintas-buku, sistem menggunakan fungsi
    \texttt{infer\_cross\_links} sebagai cadangan. Fungsi ini mencari pasangan
    konsep antarbuku yang banyak menggunakan kata yang sama (serupa dengan cara
    kerja Jaccard), namun hanya menghitung kata sepanjang minimal empat huruf dan
    bukan \textit{stopword}, lalu mengambil dua pasangan dengan kemiripan
    tertinggi. Berbeda dengan dua langkah lainnya, langkah ini menambah tautan,
    bukan membuangnya.

    \item \textit{Membersihkan relasi.} Terakhir, \textit{triple} yang tipe
    relasi atau targetnya kosong ikut dibuang agar tidak ada relasi cacat yang
    masuk ke graf.
  \end{enumerate}
\end{enumerate}

% ============================================================
\section{Graph Construction (Eksekusi Neo4j Cypher)}\label{sec:graph-construction}
% ============================================================

Tahap ini membangun \textit{knowledge graph} ke Neo4j dengan strategi
tiga-\textit{pass}, yaitu tiga tahap yang dijalankan berurutan. Urutan ini
dipilih agar graf yang terbentuk konsisten secara referensial dan idempoten. Maksud konsisten secara referensial,
setiap relasi selalu menunjuk ke \textit{node} yang sudah ada, sehingga tidak
muncul relasi yang menggantung ke konsep yang belum sempat dibuat. Karena itu,
seluruh \textit{node} dibentuk lebih dahulu, baru relasinya ditarik kemudian.
Adapun idempoten, proses ini bisa dijalankan ulang berkali-kali tanpa membuat
\textit{node} atau relasi ganda, sebab semua operasi pembuatan \textit{node} dan
relasi memakai \texttt{MERGE} yang hanya membuat data bila belum ada, bukan
\texttt{CREATE} yang selalu menambah data baru.

\begin{enumerate}
  \item \textbf{Pass 1 --- Instansiasi \textit{node} dan kerangka hierarki.}
  \begin{itemize}
    \item Node \texttt{Grade}, \texttt{Chapter}, \texttt{Subtopic}, dan
      \texttt{Concept} dibuat melalui operasi \texttt{MERGE}
      (\textit{idempotent upsert}).
    \item Properti \textit{node} ditetapkan melalui klausa \texttt{ON CREATE} dan
      \texttt{ON MATCH}.
    \item Kunci \texttt{MERGE} gabungan (\texttt{name}, \texttt{grade}) menegakkan
      sifat \textit{shared \textit{node}} pada \texttt{Concept} sebagaimana didefinisikan
      pada Subbab~\ref{subsec:skema-ontologi}.
    \item Pada \textit{pass} yang sama, kerangka hierarki graf dibentuk, yaitu
      rangka struktural berjenjang dari Kelas, Bab, Subbab, hingga Konsep melalui
      relasi struktural \texttt{HAS\_CHAPTER}, \texttt{HAS\_SUBTOPIC},
      \texttt{HAS\_CONCEPT}, dan \texttt{NEXT\_CHAPTER} (urutan antarbab). Berbeda
      dengan relasi semantik yang diprediksi LLM, relasi struktural ini ditarik
      langsung dari struktur daftar isi sehingga keterurutan dan susunan
      hierarkinya terjamin benar.
  \end{itemize}

  \item \textbf{Pass 2 --- Relasi dalam-buku.}
  \begin{itemize}
    \item Relasi semantik antar-\texttt{Concept} ditanam dengan dua cara,
      bergantung pada apakah konsep tujuannya sudah ada atau belum. Apabila
      konsep tujuan sudah terdaftar sebagai \texttt{Concept} hasil ekstraksi
      pada buku yang sedang diproses, sistem langsung membuat \textit{edge}
      antar-\texttt{Concept} yang sebenarnya. Sebaliknya, apabila konsep tujuan
      belum ada, sistem tidak memaksakan pembuatan \texttt{Concept} baru,
      melainkan merepresentasikannya sebagai \textit{node} \texttt{ConceptTarget}
      yang berperan sebagai penampung sementara (\textit{placeholder}).
    \item Pembedaan label \texttt{Concept} dan \texttt{ConceptTarget}
      (Subbab~\ref{subsec:skema-ontologi}) berguna untuk audit. Label
      \texttt{Concept} menandai konsep yang benar-benar diekstrak dari teks buku
      (memiliki dasar tekstual), sedangkan \texttt{ConceptTarget} menandai konsep
      yang baru disebut sebagai tujuan relasi tetapi tidak terdapat dalam
      himpunan konsep terekstrak resmi pada bab tersebut. Karena keduanya
      berlabel berbeda, konsep yang belum memiliki
      dasar tekstual dapat ditelusuri dengan mudah untuk diperiksa lebih lanjut.
    \item Tipe relasi dinormalisasi (karakter non-alfanumerik diganti garis
      bawah) agar sesuai dengan sintaksis \textit{relationship type} Cypher.
  \end{itemize}

  \item \textbf{Pass 3 --- Relasi lintas-buku.}
  \begin{itemize}
    \item \textit{Edge} lintas-buku diberi prefiks \texttt{LINTAS\_BUKU\_}
      diikuti tipe relasinya (mis.\ \texttt{LINTAS\_BUKU\_APLIKASI\_DARI}).
    \item Diterapkan aturan pemeriksaan keabsahan agar setiap relasi lintas-buku
      benar-benar valid, yaitu buku tujuan harus berbeda dari buku sumber (supaya
      relasinya memang menyeberang antarbuku, bukan di dalam buku yang sama) dan
      buku tujuan harus ada di dalam korpus yang diproses (supaya relasinya tidak
      mengarah ke buku yang tidak tersedia).
    \item Apabila konsep tujuan belum ada pada lingkup mata pelajaran tujuan,
      \textit{node} \texttt{Concept} untuk konsep tersebut baru dibuat tepat saat
      dibutuhkan, yaitu ketika relasi lintas-buku hendak dibentuk. Cara ini
      memastikan relasi selalu memiliki \textit{node} tujuan yang valid tanpa
      perlu menyiapkan seluruh konsep di awal.
  \end{itemize}

  \item \textbf{Sifat idempoten.} Seluruh \textit{pass} menggunakan
  \texttt{MERGE} sehingga eksekusi ulang \textit{notebook} tidak menggandakan
  \textit{node} maupun \textit{edge}. Properti ini penting agar eksperimen dapat
  dijalankan secara iteratif tanpa mengotori graf yang telah terbentuk. Kueri
  Cypher untuk ketiga \textit{pass} konstruksi beserta \textit{constraint}
  keunikannya disajikan pada \ref{lamp:cypher-konstruksi}.
\end{enumerate}

% ============================================================
\section{Graph Completion (Vector Index ANN \& LLM Classifier)}\label{sec:graph-completion}
% ============================================================

Tahap ini menemukan relasi implisit lintas-disiplin yang tidak ditangkap LLM
pada Tahap 2, mengikuti paradigma \textit{embedding} + \textit{retrieval} + \textit{LLM-as-classifier}.
Prosesnya terdiri dari empat sub-tahap: (1) \textit{concept embedding},
(2) \textit{candidate retrieval} via ANN, (3) \textit{candidate pruning},
dan (4) \textit{relation classification} via LLM.

Perlu ditegaskan bahwa tahap ini tetap tergolong \textit{Knowledge Graph
Completion} (KGC), bukan sekadar penemuan relasi (\textit{relation discovery})
berbasis teks. \textcite{ji_survey_2022} mendefinisikan KGC sebagai tugas
menambahkan \textit{triple} baru ke dalam graf yang sudah ada, dengan subtugas
mencakup \textit{link prediction}, \textit{entity prediction}, dan
\textit{relation prediction}. Tahap \textit{completion} dalam penelitian ini
beroperasi di atas himpunan \texttt{Concept} yang telah terbentuk pada tahap
konstruksi, lalu memprediksi relasi lintas-disiplin implisit yang belum tertangkap
pada tahap ekstraksi (Tahap~2), di antara konsep-konsep yang sudah ada; dengan
demikian prosesnya merupakan \textit{relation prediction}, yaitu
salah satu subtugas KGC, dan berbeda dari \textit{relation extraction} atas teks
mentah yang justru dikerjakan pada tahap ekstraksi (Tahap~2). Penggunaan
\textit{retrieval} berbasis \textit{embedding} untuk menyaring kandidat dan LLM
untuk mengklasifikasikan relasi sejalan dengan paradigma KGC berbasis LLM:
\textcite{yao_exploring_2025} memperlakukan KGC sebagai tugas tekstual dengan LLM
untuk \textit{triple classification} dan \textit{relation prediction};
\textcite{wei_kicgpt_2023} memadukan komponen pencarian kandidat (\textit{retriever})
dengan LLM untuk KGC;
serta \textcite{pan_unifying_2024} menggolongkan pemanfaatan LLM untuk melengkapi
relasi pada KG seperti ini sebagai salah satu kelas pendekatan
\textit{LLM-augmented} KGC. Dengan demikian, langkah \textit{retrieval} berperan sebagai
mekanisme skalabilitas untuk membatasi ruang kandidat, bukan penyimpangan dari
tugas KGC.

\begin{enumerate}
  \item \textbf{\texttt{Concept} embedding.}

  Atribut \texttt{description} pada tiap \texttt{Concept} di-\textit{encode}
  menjadi vektor padat (\textit{dense vector}):
  \[
    \mathbf{e}_c \in \mathbb{R}^{d}, \quad d = 3072
  \]
  menggunakan model \textit{embedding} \texttt{gemini-embedding-001} melalui Google GenAI SDK.
  Representasi vektor ini memetakan semantik teks ke ruang Euclidean berdimensi
  tinggi sehingga konsep yang secara makna berkaitan akan memiliki vektor yang
  berdekatan. Vektor yang dihasilkan disimpan sebagai properti \textit{node} (embedding) di
  Neo4j dan diindeks menggunakan Neo4j Vector Index. Perlu ditegaskan bahwa
  \textit{embedding} ini bersifat semantik (diturunkan dari deskripsi tekstual
  konsep) sehingga berbeda dari \textit{embedding} struktural berbasis topologi
  graf yang keterbatasannya dibahas pada Subbab~\ref{subsec:kgc}. Selain itu,
  \textit{embedding} di sini hanya berperan sebagai penyaring kandidat, sedangkan
  keputusan akhir relasi tetap ditentukan oleh LLM.

  Pengindeksan memanfaatkan algoritma HNSW (lihat Subbab~\ref{subsec:kgc}), yang
  struktur graf berlapisnya diilustrasikan pada Gambar~\ref{fig:hnsw}. Parameter
  HNSW tidak dikonfigurasi secara eksplisit pada \texttt{CREATE VECTOR INDEX},
  sehingga Neo4j menerapkan nilai \textit{default} ($M = 16$,
  \texttt{ef\_construction} $= 100$) dengan \textit{vector quantization} aktif
  untuk mengoptimalkan penggunaan memori.

  \item \textbf{Candidate retrieval via ANN (Approximate Nearest Neighbor).}

  Untuk tiap \textit{query} \textit{node} $q$, dijalankan pencarian ANN pada Neo4j Vector Index
  yang mengembalikan top-$k$ \textit{node} terdekat berdasarkan \textit{cosine similarity}
  (didefinisikan pada Subbab~\ref{subsec:kgc}). Proyeksi 2D
  (UMAP) atas seluruh konsep antar mapel ditunjukkan pada
  Gambar~\ref{fig:concept-umap}: konsep diwarnai berdasarkan mata
  pelajaran sehingga pengelompokan intra-mapel tampak jelas, sementara
  pasangan lintas-buku yang lolos ambang $\tau$ ditandai sebagai \textit{edge}
  ANN.

  \begin{figure}[htbp]
    \centering
    \includegraphics[width=\maxwidth]{assets/pics/concept-umap.png}
    \caption{Proyeksi 2D (UMAP) sebaran konsep antar mapel pada indeks
      vektor; titik diwarnai per mata pelajaran, \textit{edge} ANN ditampilkan
      untuk pasangan lintas-buku dengan \textit{cosine similarity} $\geq \tau$
      ($\tau = 0{,}70$).}
    \label{fig:concept-umap}
  \end{figure}

  Hanya pasangan dengan $\text{sim} \geq \tau$ (dengan $\tau = 0{,}70$ untuk
  pencarian lintas-buku) yang lolos ke tahap berikutnya:
  \[
    \text{candidates}(q)
    = \bigl\{\, c \in \mathcal{C} \;\big|\; \text{sim}(\mathbf{e}_q,\, \mathbf{e}_c) \geq \tau \,\bigr\}
  \]
  Berkat indeks HNSW (Subbab~\ref{subsec:kgc}), pencarian kandidat menghindari
  perbandingan \textit{pairwise} $O(n^2)$ sehingga tetap \textit{tractable} pada
  skala ratusan hingga ribuan konsep.

  \item \textbf{Candidate pruning.}

  Pasangan kandidat yang lolos ANN disaring melalui empat filter sebelum
  diklasifikasi LLM:
  \begin{enumerate}[label=\roman*.]
    \item \textit{Eksklusi self-pair}: kueri ANN mencegah sebuah \textit{node}
      menjadi kandidat bagi dirinya sendiri.
    \item \textit{Deduplikasi}: setiap pasangan dinormalisasi menjadi urutan
      tetap sehingga $(A,B)$ dan $(B,A)$ diperlakukan identik dan hanya diproses
      sekali.
    \item \textit{Penyaringan lintas-mapel}: pasangan dari mata pelajaran yang
      sama dibuang karena relasi dalam-mapel telah ditangani pada Tahap~2,
      sehingga tahap \textit{completion} hanya mencari jembatan antardisiplin.
    \item \textit{Penyaringan relasi tertipe}: pasangan yang sudah memiliki
      \textit{edge} bertipe (selain \texttt{SIMILAR\_TO}) dibuang untuk mencegah
      duplikasi relasi.
  \end{enumerate}
  Tahap \textit{pruning} ini secara signifikan menekan jumlah pasangan yang
  diumpankan ke LLM, sehingga mengurangi biaya inferensi dan menjaga
  \textit{precision}.

  \item \textbf{Relation classification via LLM.}

  Tiap pasangan yang lolos \textit{pruning} diklasifikasikan oleh LLM
  (Gemini 2.5 Flash Lite) melalui \textit{prompt few-shot} yang memuat: (1)
  deskripsi dan konteks kedua konsep (nama, deskripsi, mata pelajaran); (2)
  konteks graf berupa bab asal, hingga lima \textit{edge} keluar yang sudah ada,
  serta tetangga bersama (\textit{shared neighbors}); dan (3) definisi ketat tiap
  label target. Untuk klasifikasi lintas-buku, LLM memilih salah satu dari lima
  label semantik \texttt{LINTAS\_BUKU\_*} yang didefinisikan pada
  Subbab~\ref{subsec:skema-ontologi}, atau label \texttt{none} apabila kemiripan
  semantik hanya superfisial sehingga tidak ada relasi kurikuler bermakna.
  Klasifikasi dilengkapi skor \textit{confidence} dalam rentang $[0, 1]$. Hasil
  klasifikasi di-\textit{cache} dengan kunci yang tidak bergantung pada urutan
  pasangan dan menyertakan \textit{sidik jari} (\textit{hash}) konteks graf,
  sehingga perubahan graf otomatis memicu klasifikasi ulang. Templat
  \textit{system prompt} dan \textit{user prompt} klasifikasi relasi lintas-buku
  beserta skema respons JSON-nya diberikan pada
  \ref{lamp:prompt-completion}.


  \item \textbf{Edge writeback.}

  Hanya pasangan dengan label non-none yang dituliskan kembali sebagai \textit{edge}
  baru ke Neo4j beserta properti \textit{confidence score}, description (satu kalimat
  \textit{rationale} Bahasa Indonesia), dan method (identifier model dan versi
  \textit{prompt}). Strategi ini melengkapi graf tanpa menambah \textit{noise} terstruktur,
  menghasilkan \textit{knowledge graph} yang lebih lengkap secara topologis sambil
  menjaga kualitas relasi melalui \textit{gating} LLM.
\end{enumerate}

% ============================================================
\section{Arsitektur Instrumen Evaluasi}\label{sec:instrumen-evaluasi}
% ============================================================

Untuk mengukur efektivitas dan kualitas \textit{Knowledge Graph} (KG) yang telah
dibangun, penelitian ini mengembangkan instrumen evaluasi khusus yang
terdiri dari dua komponen arsitektur utama: kalkulasi metrik struktural
berbasis \textit{Cypher} dan aplikasi web KG Review App untuk validasi pakar.

\begin{enumerate}
  \item \textbf{Implementasi Kalkulasi Struktural pada Neo4j Aura.}

  Seluruh perhitungan metrik topologi graf dieksekusi secara langsung di sisi
  basis data menggunakan kueri \textit{Cypher} pada Neo4j Aura. Pendekatan
  eksekusi \textit{database-side} ini memastikan bahwa metrik dapat diturunkan ulang
  setiap saat melalui kueri \textit{read-only}. Pendekatan ini dirancang khusus agar
  kompatibel dengan spesifikasi tier Aura Free yang tidak menyertakan pustaka
  \textit{Graph Data Science} (GDS), sehingga hasil perhitungannya dapat langsung
  diintegrasikan ke dalam \textit{dashboard} visualisasi.

  \item \textbf{Arsitektur KG Review App.}

  Untuk memfasilitasi validasi oleh pakar (\textit{human evaluation}),
  dikembangkan aplikasi web kustom bernama KG Review App, sebuah aplikasi
  berbasis Next.js dengan penyimpanan Redis. Dua fitur instrumen yang menjamin
  kualitas pengumpulan data ditekankan di sini. \emph{Pertama}, \textit{auto-save}
  otomatis pada setiap aksi penilaian sehingga sesi panjang (lebih dari 100
  \textit{triple} per mata pelajaran) dapat dilanjutkan tanpa kehilangan data.
  \emph{Kedua}, penyajian \textit{triple} secara \textit{score-blind} dan
  acak (skor kepercayaan model disembunyikan dan urutan diacak) agar penilaian
  pakar tidak ter-\textit{anchor} pada keyakinan model. Rincian \textit{tech stack},
  skema kunci Redis, dan diagram arsitektur lengkap disajikan di
  \ref{lamp:kg-review-app}; tampilan antarmukanya ditunjukkan pada Gambar~\ref{fig:admin-dashboard}.

  \begin{figure}[htbp]
    \centering
    \includegraphics[width=\maxwidth]{assets/pics/admin-dashboard.png}
    \caption{Tampilan dasbor Admin untuk mengunggah graf JSON pada KG Review App.}
    \label{fig:admin-dashboard}
  \end{figure}

  \item \textbf{Role-Based Access Control (RBAC).}

  Untuk menjaga integritas proses validasi, aplikasi menerapkan kontrol akses
  berbasis peran dengan dua peran utama: \textbf{Admin}, yang memegang hak
  \textit{management} (mengunggah graf, mengatur publikasi, serta meninjau metrik respons
  pakar secara \textit{read-only}); dan \textbf{Expert} (\textit{pakar}), yang
  memegang hak \textit{read-write} untuk menilai relasi dan mengusulkan
  \textit{missing triples}. Prosedur penilaian per peran telah dirinci pada
  Subbab~\ref{subsec:cara-pengumpulan}.

  \item \textbf{Desain Skema Validasi Dua Fase.}

  Antarmuka evaluasi untuk Expert mengikuti dua fase validasi yang telah
  didefinisikan pada Subbab~\ref{subsec:kerangka-konsep} (Fase~1 intra-mapel dan
  Fase~2 antar-mapel). Pemisahan ini diterapkan pada tingkat antarmuka untuk
  menurunkan beban kognitif (\textit{cognitive load}) pakar: pakar menilai
  relasi dalam satu mata pelajaran terlebih dahulu sebelum mengevaluasi relasi
  lintas-disiplin, sehingga setiap fase dapat dianalisis secara lebih terisolasi
  dan akurat.

  \item \textbf{Pipeline Konsolidasi Konsensus (\textit{LLM-as-a-judge}).}

  Komponen terakhir mengubah penilaian mentah pakar menjadi \textit{ground truth}
  final dan graf konsensus. Prosedur konseptualnya telah diuraikan pada
  Subbab~\ref{sec:pengolahan-data}; bagian ini memaparkan implementasinya
  melalui tiga \textit{notebook} yang dijalankan berurutan: pembentukan
  konsensus, perhitungan metrik, dan pemuatan graf konsensus ke Neo4j.

  Masukannya adalah dua berkas penilaian pakar per mata pelajaran yang diekspor
  dari KG Review App (Fase~1) beserta berkas penilaian lintas-mapel (Fase~2).
  \textit{Notebook} konsolidasi mengagregasi kedua penilaian; bila kedua pakar
  berbeda rating atas satu \textit{triple}, atau bila sebuah \textit{missing
  triple} hanya diusulkan satu pakar, kasus tersebut diteruskan ke
  \textit{LLM-as-a-judge} (Gemini~3.1 Pro Preview) sebagaimana kerangka
  pada Gambar~\ref{fig:llm-judge}. Agar putusan AI Judge berlandaskan sumber
  alih-alih pengetahuan parametrik model, \textit{prompt} penengahan disertai
  potongan konteks buku teks yang relevan (hasil ekstraksi Tahap~1). Setiap
  panggilan AI Judge disimpan sebagai \textit{checkpoint} sehingga eksekusi ulang
  tidak memanggil API secara berulang, sementara putusannya tetap dapat
  ditelusuri (\textit{auditable}) dan direproduksi.

  Keluarannya ada dua. \emph{Pertama}, \textit{gold standard} per mata pelajaran
  berupa daftar label final yang menjadi masukan perhitungan metrik
  (Gwet's AC1 serta \textit{Fuzzy Precision}/\textit{Recall}/F1). \emph{Kedua},
  berkas KG konsensus, yakni ekspor graf dengan label konsensus dan metadata
  peninjauan yang ditanamkan pada tiap relasi. \textit{Notebook} 
  menulis KG konsensus ini kembali ke Neo4j dengan \textit{constraint} keunikan,
  membentuk \textit{node} \texttt{Grade}, \texttt{Chapter}, \texttt{Subtopic},
  \texttt{Concept}, dan \texttt{ConceptTarget} beserta relasi struktural dan relasi
  semantik konsensus. Untuk analitik topologi, perhitungan
  metrik (Subbab~\ref{subsec:struktural}) berpijak pada relasi semantik
  antar-\texttt{Concept}, bukan relasi struktural hierarkis, agar nilai
  \textit{modularity} dan komunitas tidak terdistorsi oleh struktur hierarki
  bab. Volume dan hasil penengahan AI Judge atas data ini dilaporkan pada
  Subbab~\ref{subsec:validasi-pakar} (Tabel~\ref{tab:llm-judge}).
\end{enumerate}

\ifSubfilesClassLoaded{\printbibliography[title={Daftar Pustaka}]}{}

\end{document}
